% GENERATED FILE. Edit canonical lessons, then run npm run generate:books.
% canonical-source-hash: fnv1a64:fc8b9194e79bbce3
% canonical-lessons: JA-C349-yosou, JA-C349-joukyou, JA-C349-kikai, JA-C349-nogasu, JA-C349-sekkaku

\chapter{Expecting the Unexpected}
\label{ch:ja-expecting-the-unexpected}
\hlchaptermodality{349}

\begin{chapteropening}
\textbf{By the end of this chapter:} \emph{I can say what we will do if things turn out differently.}

Complete the last lesson of chapter 349: I can say what we will do if things turn out differently.
\end{chapteropening}

\section[\texorpdfstring{\ja{よそう} (yosō)}{よそう (yosō)}]{\ja{よそう} (yosō) — an expectation; a forecast}

\label{lesson:JA-C349-yosou}



\begin{quote}
Before the new one: say the Japanese for the hour, then the Japanese for consecutive holidays.
\end{quote}

\subsection*{You'll want to know: \ja{よそう}}
\textbf{\ja{よそう}} — \emph{yosō} — \textquotedblleft{}an expectation; a forecast\textquotedblright{}.

\begin{quote}
\textbf{\ja{けっかが} \ja{よそうと} \ja{ちがったら、もう} \ja{いちど} \ja{しらべます。}} — \emph{kekka ga yosō to chigattara, mō ichido shirabemasu} — \textquotedblleft{}If the result is not what we expected, we will investigate again.\textquotedblright{}
\end{quote}

\begin{cousinweb}
\textbf{\ja{よそう}} (\emph{yosō}) is Sino-Japanese (kango): two Chinese characters read with their on'yomi, readings borrowed from Chinese. They are \emph{yo}, \textquotedblleft{}beforehand\textquotedblright{} (Middle Chinese \emph{yoH}) and \emph{sō}, \textquotedblleft{}idea\textquotedblright{} (\emph{sjangX}) (KANJIDIC2; Middle Chinese in Baxter's transcription).
\end{cousinweb}

\begin{grammarlens}[title={What you've built}]
One more word for saying what happens if.
\end{grammarlens}

\subsection*{Your turn}
\begin{itemize}
\raggedright
  \item \emph{Say it:} \emph{yosō}
  \item \emph{Say it:} \emph{yosō}, once more
  \item \emph{Say it:} \emph{kekka ga yosō to chigattara, mō ichido shirabemasu} — If the result is not what we expected, we will investigate again.
  \item \emph{Recall:} say the Japanese for the hour, then the Japanese for consecutive holidays, then say \emph{yosō} again
\end{itemize}

\begin{tcolorbox}[breakable,colback=teal!4,colframe=teal!35!black,title={Before you move on}]
What does \ja{よそう} mean? (\textquotedblleft{}An expectation; a forecast\textquotedblright{}.) Say it once more.
\end{tcolorbox}
\section[\texorpdfstring{\ja{じょうきょう} (jōkyō)}{じょうきょう (jōkyō)}]{\ja{じょうきょう} (jōkyō) — the situation; the state of things}

\label{lesson:JA-C349-joukyou}



\begin{quote}
Before the new one: say the Japanese for consecutive holidays, then the Japanese for an expectation.
\end{quote}

\subsection*{You'll want to know: \ja{じょうきょう}}
\textbf{\ja{じょうきょう}} — \emph{jōkyō} — \textquotedblleft{}the situation; the state of things\textquotedblright{}.

\begin{quote}
\textbf{\ja{じょうきょうが} \ja{かわったら、れんらくします。}} — \emph{jōkyō ga kawattara, renraku shimasu} — \textquotedblleft{}If the situation changes, I will get in touch.\textquotedblright{}
\end{quote}

\begin{cousinweb}
\textbf{\ja{じょうきょう}} (\emph{jōkyō}) is Sino-Japanese (kango): two Chinese characters read with their on'yomi, readings borrowed from Chinese. They are \emph{jō}, \textquotedblleft{}conditions\textquotedblright{} (Middle Chinese \emph{dzrjangH}) and \emph{kyō}, \textquotedblleft{}condition\textquotedblright{} (\emph{xjwangH}) (KANJIDIC2; Middle Chinese in Baxter's transcription).
\end{cousinweb}

\begin{grammarlens}[title={What you've built}]
One more word for saying what happens if.
\end{grammarlens}

\subsection*{Your turn}
\begin{itemize}
\raggedright
  \item \emph{Say it:} \emph{jōkyō}
  \item \emph{Say it:} \emph{jōkyō}, once more
  \item \emph{Say it:} \emph{jōkyō ga kawattara, renraku shimasu} — If the situation changes, I will get in touch.
  \item \emph{Recall:} say the Japanese for consecutive holidays, then the Japanese for an expectation, then say \emph{jōkyō} again
\end{itemize}

\begin{tcolorbox}[breakable,colback=teal!4,colframe=teal!35!black,title={Before you move on}]
What does \ja{じょうきょう} mean? (\textquotedblleft{}The situation; the state of things\textquotedblright{}.) Say it once more.
\end{tcolorbox}
\section[\texorpdfstring{\ja{きかい} (kikai)}{きかい (kikai)}]{\ja{きかい} (kikai) — an opportunity; a chance}

\label{lesson:JA-C349-kikai}



\begin{quote}
Before the new one: say the Japanese for an expectation, then the Japanese for the situation.
\end{quote}

\subsection*{You'll want to know: \ja{きかい}}
\textbf{\ja{きかい}} — \emph{kikai} — \textquotedblleft{}an opportunity; a chance\textquotedblright{}.

\begin{quote}
\textbf{\ja{きかいが} \ja{あったら、}\ja{日本}\ja{に} \ja{いきたいです。}} — \emph{kikai ga attara, nihon ni ikitai desu} — \textquotedblleft{}If I get the chance, I want to go to Japan.\textquotedblright{}
\end{quote}

\begin{cousinweb}
\textbf{\ja{きかい}} (\emph{kikai}) is Sino-Japanese (kango): two Chinese characters read with their on'yomi, readings borrowed from Chinese. They are \emph{ki}, \textquotedblleft{}opportunity\textquotedblright{} (Middle Chinese \emph{kj+j}) and \emph{kai}, \textquotedblleft{}meeting\textquotedblright{} (\emph{hwajH}) (KANJIDIC2; Middle Chinese in Baxter's transcription).
\end{cousinweb}

\begin{grammarlens}[title={What you've built}]
One more word for saying what happens if.
\end{grammarlens}

\subsection*{Your turn}
\begin{itemize}
\raggedright
  \item \emph{Say it:} \emph{kikai}
  \item \emph{Say it:} \emph{kikai}, once more
  \item \emph{Say it:} \emph{kikai ga attara, nihon ni ikitai desu} — If I get the chance, I want to go to Japan.
  \item \emph{Recall:} say the Japanese for an expectation, then the Japanese for the situation, then say \emph{kikai} again
\end{itemize}

\begin{tcolorbox}[breakable,colback=teal!4,colframe=teal!35!black,title={Before you move on}]
What does \ja{きかい} mean? (\textquotedblleft{}An opportunity; a chance\textquotedblright{}.) Say it once more.
\end{tcolorbox}
\section[\texorpdfstring{\ja{のがす} (nogasu)}{のがす (nogasu)}]{\ja{のがす} (nogasu) — to miss a chance; to let slip}

\label{lesson:JA-C349-nogasu}



\begin{quote}
Before the new one: say the Japanese for the situation, then the Japanese for an opportunity.
\end{quote}

\subsection*{You'll want to know: \ja{のがす}}
\textbf{\ja{のがす}} — \emph{nogasu} — \textquotedblleft{}to miss a chance; to let slip\textquotedblright{}.

\begin{quote}
\textbf{\ja{この} \ja{きかいを} \ja{のがしたら、つぎは} \ja{らいねんです。}} — \emph{kono kikai o nogashitara, tsugi wa rainen desu} — \textquotedblleft{}If we miss this chance, the next one is next year.\textquotedblright{}
\end{quote}

Polite \textbf{\ja{のがします}}, \emph{nogashimasu}; a godan verb. Past (polite): \textbf{\ja{あには} \ja{きかいを} \ja{のがしました。}}, \emph{ani wa kikai o nogashimashita}.

\begin{cousinweb}
\textbf{\ja{のがす}} (\emph{nogasu}) is a native Japanese word (wago), inherited rather than borrowed from Chinese, and its older origin is not settled. Its kanji was attached to it later as a kun'yomi, a native reading given to a Chinese character (JMdict, KANJIDIC2).
\end{cousinweb}

\begin{grammarlens}[title={What you've built}]
One more word for saying what happens if.
\end{grammarlens}

\subsection*{Your turn}
\begin{itemize}
\raggedright
  \item \emph{Say it:} \emph{nogasu}
  \item \emph{Say it:} \emph{nogasu}, once more
  \item \emph{Say it:} \emph{kono kikai o nogashitara, tsugi wa rainen desu} — If we miss this chance, the next one is next year.
  \item \emph{Recall:} say the Japanese for the situation, then the Japanese for an opportunity, then say \emph{nogasu} again
\end{itemize}

\begin{tcolorbox}[breakable,colback=teal!4,colframe=teal!35!black,title={Before you move on}]
What does \ja{のがす} mean? (\textquotedblleft{}To miss a chance; to let slip\textquotedblright{}.) Say it once more.
\end{tcolorbox}
\section[\texorpdfstring{\ja{せっかく} (sekkaku)}{せっかく (sekkaku)}]{\ja{せっかく} (sekkaku) — at great pains; specially}

\label{lesson:JA-C349-sekkaku}



\begin{quote}
Before the new one: say the Japanese for an opportunity, then the Japanese for to miss a chance.
\end{quote}

\subsection*{You'll want to know: \ja{せっかく}}
\textbf{\ja{せっかく}} — \emph{sekkaku} — \textquotedblleft{}at great pains; specially\textquotedblright{}.

\begin{quote}
\textbf{\ja{もし} \ja{あめが} \ja{ふったら、せっかくの} \ja{りょこうが} \ja{ちゅうしに} \ja{なります。}} — \emph{moshi ame ga futtara, sekkaku no ryokō ga chūshi ni narimasu} — \textquotedblleft{}If it rains, the trip we planned so carefully will be called off.\textquotedblright{}
\end{quote}

\begin{cousinweb}
\textbf{\ja{せっかく}} (\emph{sekkaku}) is Sino-Japanese (kango): two Chinese characters read with their on'yomi, readings borrowed from Chinese. They are \emph{setsu}, \textquotedblleft{}fold\textquotedblright{} (Middle Chinese \emph{tsyet}) and \emph{kaku}, \textquotedblleft{}angle\textquotedblright{} (\emph{kaewk}) (KANJIDIC2; Middle Chinese in Baxter's transcription). The two join as \emph{sek-kaku}; JMdict notes that it is usually written in kana.
\end{cousinweb}

\begin{grammarlens}[title={What you've built}]
One more word for saying what happens if.
\end{grammarlens}

\subsection*{Your turn}
\begin{itemize}
\raggedright
  \item \emph{Say it:} \emph{sekkaku}
  \item \emph{Say it:} \emph{sekkaku}, once more
  \item \emph{Say it:} \emph{moshi ame ga futtara, sekkaku no ryokō ga chūshi ni narimasu} — If it rains, the trip we planned so carefully will be called off.
  \item \emph{Recall:} say the Japanese for an opportunity, then the Japanese for to miss a chance, then say \emph{sekkaku} again
\end{itemize}

\begin{tcolorbox}[breakable,colback=teal!4,colframe=teal!35!black,title={Before you move on}]
What does \ja{せっかく} mean? (\textquotedblleft{}At great pains; specially\textquotedblright{}.) Say it once more.
\end{tcolorbox}
